\documentclass[10pt,a4paper]{amsart}
\usepackage[margin=19mm]{geometry}
\usepackage{amsmath,amssymb,mathtools}
\usepackage[scr=rsfs,cal=euler]{mathalfa}
\usepackage{xfrac}
\usepackage[unicode,hidelinks,bookmarks=false]{hyperref}
\newcommand{\ie}{i.e.\ }
\newcommand{\cf}{cf.\ }
\newcommand{\resp}{respectively\ }
\newcommand{\Subsection}{\subsection}
\providecommand{\nobreakdash}{\nobreak\hbox{-}}
\setlength{\emergencystretch}{2em}
\multlinegap0pt
\usepackage{bm}
\usepackage{bbm}
\usepackage{dsfont}
\usepackage{xspace}
\usepackage{cancel}
\usepackage{mathtools}
\usepackage{amsthm}
\makeatletter
\@ifundefined{theorem}{\newtheorem{theorem}{Theorem}}{}
\@ifundefined{proposition}{\newtheorem{proposition}[theorem]{Proposition}}{}
\makeatother
\title[A survey on the Asplund property for spaces $C(X)$ of contin]{A survey on the Asplund property for spaces $C(X)$ of continuous functions}
\author{Marian Fabian, Jerzy Kakol, Arkady Leiderman}
\date{Source: arXiv:2609.23167v1 (CC BY 4.0)}
\begin{document}
\maketitle

\noindent\emph{Source and licence.} A survey on the Asplund property for spaces $C(X)$ of continuous functions. Version arXiv:2609.23167v1 (CC BY 4.0), distributed under the Creative Commons Attribution 4.0 International licence, as recorded in the arXiv licence metadata for this exact version and shown on the abstract page https://arxiv.org/abs/2609.23167v1. Full rights notice: licenses/arxiv-2609.23167-CC-BY-4.0.txt.\par\smallskip

\noindent\textit{Editorial scope.} Counted unit: the Proposition whose source label is \texttt{None} in the article (source0.tex lines 1067--1086, source label \texttt{None}), delivered together with the article's own complete original proof of it (source0.tex lines 1087--1334, one \texttt{proof} environment of 248 lines). No mathematics is written, repaired or re-derived: statement and proof are the article's own text line for line. Required context retained uncounted: none. Every definition, display and label of those lines is retained unchanged. Omitted from this package and not redistributed: 1--1066, 1335--1825. Nothing inside a retained excerpt is omitted or altered: every retained body is delivered exactly as the article's lines are written, so no omission receipt of any kind accompanies this package. Citations: the retained excerpts carry no citation command at all, so no bibliography entry of the article is transcribed and no reference is redistributed. Reference adaptation: none was needed and none was made; every reference inside the delivered unit resolves inside it, so no unresolved reference or citation is produced. Printed slips kept literal: no printed word, symbol, hypothesis or number of the counted statement or of its proof was altered. Layout and class: the article's own class is \texttt{amsart} with packages including amsmath, amssymb, eucal, xy, latexsym, amstext, amsfonts, amsthm, amsopn, amsbsy, layout, color, graphicx, bm; the wrapper is the lane's pinned \texttt{amsart} editorial wrapper at 10pt on A4, which restates the theorem-like environments the retained text needs and adds the article's own macro definitions that the retained text needs (none), with extra wrapper packages (bm, bbm, dsfont, xspace, cancel, mathtools). The delivered numbering is the unit's own. Assets: no retained span contains a figure environment, a graphics-inclusion command, a PDF-page inclusion, a table or a TikZ picture; no image file is copied into this package, the package declares no asset dependency and no image file is redistributed with it. Licence: version arXiv:2609.23167v1 of this article is distributed under the Creative Commons Attribution 4.0 International licence, as recorded in the arXiv licence metadata for this exact version and shown on the abstract page https://arxiv.org/abs/2609.23167v1. A full rights notice is delivered at licenses/arxiv-2609.23167-CC-BY-4.0.txt. Characters: every retained excerpt is pure ASCII, so the delivered engine, XeLaTeX with its default Latin Modern text font, reports no missing character.
\begin{proposition}
Let $E$ be a locally convex space. Assume that  $(A_n)$ is an increasing sequence of closed absolutely convex subsets of $E$ covering $E$
and that there exists a sequence $(x_n)$ such that
\[
x_n\in A_{n+1}\setminus \operatorname{sp}(A_n)
\qquad (n\in\mathbb N).
\]

Then the dual space $E'$ contains functionals $f_n\in A_n^\circ$ $(n\in\mathbb N)$ for which
\[
q(x):=\sup_{n\in\mathbb N}|f_n(x)|
\]
satisfies
\[
\sum_{i=1}^{\infty}|a_i|
\le
q\!\left(\sum_{i=1}^{\infty}a_i x_i\right)
\qquad\text{for all }(a_i)\in c_{00}.
\]
\end{proposition}

\begin{proof}
The idea is to construct inductively a sequence of continuous functionals
\[
f_n\in A_n^\circ,\qquad n\in\mathbb N,
\]
whose associated seminorms dominate the $\ell_1$-norm of the coefficients of finite linear combinations of the vectors $x_n$.

For every $k\in\mathbb N$ define
\[
q_k(x):=\max\{|f_1(x)|,\ldots,|f_k(x)|\}.
\]
We shall show inductively that
\begin{equation}\label{eq:induction}
\frac{k+1}{k}\sum_{i=1}^{k}|a_i|
   \le
q_k\Big(\sum_{i=1}^{k}a_ix_i\Big)
\qquad
\forall (a_i)\in c_{00}.
\end{equation}
Since $(k+1)/k\to1$, letting $k\to\infty$ will yield
\[
\sum_{i=1}^{\infty}|a_i|
   \le
q\Big(\sum_{i=1}^{\infty}a_ix_i\Big),
\]
where
\[
q(x):=\sup_{n\in\mathbb N}|f_n(x)|.
\]

\medskip

We start with the construction of $f_1$. Since
$x_1\notin 2 A_1$
and $2 A_1$ is a closed convex set, the Hahn--Banach separation theorem yields a functional
$f_1\in E'$ such that
\[
f_1(x_1)=2
   >
\sup f_1\Big[2 A_1\Big].
\]
Consequently,
$f_1\in A_1^\circ$.
Moreover,
\[
q_1(a_1x_1)=|a_1|\,|f_1(x_1)|
=2 |a_1|,
\]
and therefore \eqref{eq:induction} holds for $k=1$.

\medskip

Assume now that for some $k\in\mathbb N$ we have already constructed
\[
f_1\in A_1^\circ,\ldots,f_k\in A_k^\circ
\]
and that \eqref{eq:induction} holds.

Choose $P_0>0$ so large that
\begin{equation}\label{eq:P0}
\frac{k+2}{k+1}
<
\frac{(1+1/k)P-q_k(x_{k+1})}{P+1}
\qquad
\forall P\ge P_0.
\end{equation}

Set
\begin{equation}\label{eq:Q}
Q:=
\frac{k+2}{k+1}(P_0+1)+P_0.
\end{equation}

Since
$x_{k+1}\notin {\rm sp}(A_k)$
and
${\rm sp}(A_k)\subset QA_{k+1}$,
the Hahn--Banach theorem provides
$f_{k+1}\in E'$
such that
\begin{equation}\label{eq:HB}
f_{k+1}(x_{k+1})=Q
>
\sup f_{k+1}[QA_{k+1}].
\end{equation}
Hence
\[
f_{k+1}\in A_{k+1}^\circ.
\]

Fix
\[
x:=\sum_{i=1}^{k}a_ix_i .
\]
Since
$x_1,\ldots,x_k\in A_{k+1}$,
we obtain
\begin{equation}\label{eq:estimate21}
|f_{k+1}(\pm x)|
\le
\sum_{i=1}^{k}|a_i|.
\end{equation}

We distinguish two cases.

\medskip

\noindent
{\bf Case 1.}
\[
\sum_{i=1}^{k}|a_i|\ge P_0.
\]

Using \eqref{eq:P0}, the inductive hypothesis
\eqref{eq:induction}, and the triangle inequality, we get
\[
\frac{k+2}{k+1}
<
\frac{(1+1/k)\sum_{i=1}^{k}|a_i|-q_k(x_{k+1})}
     {\sum_{i=1}^{k}|a_i|+1}
\]
\[
\le
\frac{
q_k\!\left(\sum_{i=1}^{k}a_ix_i\right)-q_k(x_{k+1})
}{
\sum_{i=1}^{k}|a_i|+1
}
\le
\frac{
q_k(\pm x+x_{k+1})
}{
\sum_{i=1}^{k}|a_i|+1
}
\]
\[
\le
\frac{
q_{k+1}(\pm x+x_{k+1})
}{
\sum_{i=1}^{k}|a_i|+1
}.
\]
Hence
\begin{equation}\label{eq:case1}
\frac{k+2}{k+1}
\Big(
\sum_{i=1}^{k}|a_i|+1
\Big)
<
q_{k+1}(\pm x+x_{k+1}).
\end{equation}

\medskip

\noindent
{\bf Case 2.}
\[
\sum_{i=1}^{k}|a_i|<P_0.
\]

Using \eqref{eq:Q}, \eqref{eq:HB}, and
\eqref{eq:estimate21}, we obtain
\[
\frac{k+2}{k+1}
=
\frac{Q-P_0}{P_0+1}
<
\frac{
Q-\sum_{i=1}^{k}|a_i|
}{
\sum_{i=1}^{k}|a_i|+1
}
\]
\[
=
\frac{
f_{k+1}(x_{k+1})
-\sum_{i=1}^{k}|a_i|
}{
\sum_{i=1}^{k}|a_i|+1
}
\le
\frac{
f_{k+1}(x_{k+1})
-|f_{k+1}(\pm x)|
}{
\sum_{i=1}^{k}|a_i|+1
}
\]
\[
\le
\frac{
|f_{k+1}(\pm x+x_{k+1})|
}{
\sum_{i=1}^{k}|a_i|+1
}
\le
\frac{
q_{k+1}(\pm x+x_{k+1})
}{
\sum_{i=1}^{k}|a_i|+1
}.
\]
Therefore,
\begin{equation}\label{eq:case2}
\frac{k+2}{k+1}
\Big(
\sum_{i=1}^{k}|a_i|+1
\Big)
\le
q_{k+1}(\pm x+x_{k+1}).
\end{equation}

\medskip

In both cases we arrive at
\[
\frac{k+2}{k+1}
\Big(
\sum_{i=1}^{k}|a_i|+1
\Big)
\le
q_{k+1}(\pm x+x_{k+1}).
\]
Choosing the sign according to the sign of $a_{k+1}$ and replacing
$x$ by $\sum_{i=1}^{k}a_ix_i$, we obtain
\[
\frac{k+2}{k+1}
\sum_{i=1}^{k+1}|a_i|
\le
q_{k+1}
\Big(
\sum_{i=1}^{k+1}a_ix_i
\Big).
\]
Thus the inductive step is complete.

Hence \eqref{eq:induction} holds for every $k\in\mathbb N$. Passing to the limit as $k\to\infty$, we conclude that
\[
\sum_{i=1}^{\infty}|a_i|
\le
q\Big(\sum_{i=1}^{\infty}a_ix_i\Big)
\qquad
\forall (a_i)\in c_{00},
\]
which completes the proof.
\end{proof}

\end{document}
